% Part:sets-functions-relations
% Chapter: size-of-sets
% Section: equinumerous-sets

\documentclass[../../../include/open-logic-section]{subfiles}

\begin{document}

\olfileid[fr]{sfr}{siz}{equ}

\olsection{Équipotence}

Notre notion intuitive de « taille » d'un ensemble convient bien aux
ensembles finis. Mais qu'en est-il des ensembles infinis ? Pour
comparer de manière formelle les tailles de deux ensembles
\emph{quelconques}, commençons par définir ce que signifie avoir la
même taille. Voici comment Frege présente cette idée :
\begin{quote}
  Si un serveur veut s'assurer qu'il dispose sur la table autant de
  couteaux que d'assiettes, il n'a besoin de compter ni les uns ni les
  autres : il lui suffit de placer un couteau à droite de chaque
  assiette, de sorte que chaque couteau sur la table se trouve à droite
  d'une assiette. Les assiettes et les couteaux sont ainsi mis en
  correspondance de manière univoque dans les deux sens, et ce par la
  même relation de position. \citep[\S70]{Frege1884}\footnote{Traduction
  établie pour cette édition après consultation du texte allemand de
  1884 ; elle ne reproduit pas une traduction française publiée.}
\end{quote}
Une définition formelle permet de préciser l'idée de ce passage :

\begin{defn}\ollabel{comparisondef}
  On dit que $A$ est \emph{équipotent} à $B$, et on écrit
  $\cardeq{A}{B}$, si et seulement s'il existe une !!{bijection}
  $f \colon A \to B$.
\end{defn}

\begin{prop}\ollabel{equinumerosityisequi}
L'équipotence est une relation d'équivalence.\footnote{Précision
éditoriale : cela signifie ici qu'elle est réflexive, symétrique et
transitive pour les ensembles. Sur toute famille d'ensembles qui est
elle-même un ensemble, elle définit donc une relation d'équivalence
au sens du chapitre précédent ; on ne suppose pas l'existence d'un
ensemble de tous les ensembles.}
\end{prop}

\begin{proof}
Il faut montrer que l'équipotence est réflexive, symétrique et
transitive. Soient $A$, $B$ et $C$ des ensembles.

\emph{Réflexivité.} L'application identité $\Id{A} \colon A \to A$,
définie par $\Id{A}(x) = x$ pour tout $x \in A$, est une
!!{bijection}. On a donc $\cardeq{A}{A}$.

\emph{Symétrie.} Supposons que $\cardeq{A}{B}$ : il existe une
!!{bijection} $f\colon A \to B$. Puisque $f$ est !!{bijective},
sa réciproque $f^{-1}$ existe et est elle aussi !!{bijective}.
Ainsi, $f^{-1}\colon B \to A$ est une !!{bijection}, d'où
$\cardeq{B}{A}$.

\emph{Transitivité.} Supposons que $\cardeq{A}{B}$ et
$\cardeq{B}{C}$ : il existe des !!{bijection}s $f\colon A \to B$
et $g\colon B \to C$. La composée $\comp{f}{g}\colon A \to C$
est alors !!{bijective}, d'où $\cardeq{A}{C}$.
\end{proof}

\begin{prop}
Si $\cardeq{A}{B}$, alors $A$ est !!{enumerable} si et seulement si
$B$ l'est.
\end{prop}

\begin{editorial}
La démonstration suivante utilise la \olref[enm]{defn:enumerable}
si la \olref[enm]{sec} est incluse, et la
\olref[enm-alt]{defn:enumerable} dans le cas contraire.
Dans le lecteur complet, les deux versions sont présentées ci-dessous,
avec leurs définitions respectives de l'énumération.
\end{editorial}

\begin{proof}
Supposons que $\cardeq{A}{B}$, donc qu'il existe une !!{bijection}
$f \colon A \to B$, et que $A$ soit !!{enumerable}.
\frproofversions{Avec une énumération par surjection.}{Avec une énumération bijective.}{sfr:siz:enm:defn:enumerable}{
  Ou bien $A = \emptyset$, ou bien il existe une fonction
  !!{surjective} $g\colon \PosInt \to A$. Si $A = \emptyset$,
  alors $B = \emptyset$ aussi : sinon il existerait un élément
  $y \in B$ mais aucun $x \in A$ tel que $f(x) = y$,
  contrairement à la surjectivité de~$f$.\footnote{Note de traduction :
  dans les deux versions de cette démonstration, l'original écrit
  $g(x) = y$ dans l'argument relatif au cas vide. C'est la fonction
  $f\colon A \to B$ qui doit intervenir ici.}
  Si, en revanche, $g\colon \PosInt \to A$ est !!{surjective},
  alors $\comp{g}{f} \colon \PosInt \to B$ est !!{surjective}.
  En effet, soit $y \in B$. La surjectivité de $f$ fournit un
  $x \in A$ tel que $f(x) = y$ ; celle de $g$ fournit un
  $n \in \PosInt$ tel que $g(n) = x$. Par conséquent,
\[
(\comp{g}{f})(n) = f(g(n)) = f(x) = y
\]
et $\comp{g}{f}$ est bien !!{surjective}. Ainsi, $\comp{g}{f}$ est
une énumération de~$B$, donc $B$ est !!{enumerable}.}
{
  Ou bien $A = \emptyset$, ou bien il existe une !!{bijection}~$g$
  dont l'image est $A$ et dont le domaine est $\Nat$ ou un segment
  initial des entiers naturels. Si $A = \emptyset$, alors
  $B = \emptyset$ aussi : sinon il existerait un $y \in B$ sans
  aucun $x \in A$ tel que $f(x) = y$, contrairement à la
  surjectivité de~$f$.\footnote{Note de traduction : l'original
  écrit $g(x) = y$ dans le cas vide ; c'est $f\colon A \to B$
  qui doit intervenir. La même correction figure dans l'autre
  version de la démonstration.}
  Supposons donc donnée la !!{bijection}~$g$. Alors $\comp{g}{f}$
  est une !!{bijection} d'image~$B$ et de même domaine que~$g$
  (c'est-à-dire $\Nat$ ou l'un de ses segments initiaux).
  Ainsi, $B$ est !!{enumerable}.}

Si $B$ est !!{enumerable}, on montre que $A$ l'est aussi en
reprenant l'argument avec la !!{bijection} $f^{-1}\colon B \to A$
à la place de~$f$.
\end{proof}

\begin{prob}
Montrer que, si $\cardeq{A}{C}$ et $\cardeq{B}{D}$, et si
$A \cap B = C \cap D = \emptyset$, alors
$\cardeq{A \cup B}{C \cup D}$.
\end{prob}

\begin{prob}
Montrer que, si $A$ est infini et !!{enumerable}, alors
$\cardeq{A}{\Nat}$.
\end{prob}

\end{document}
